% sections/introduction.tex
% Biology-first framing; SHPN mentioned briefly as the modelling tool.
% Dense inline numbers moved to Results tables (guideline 2).
% Two heading levels only (guideline 1).

\subsection{The Problem of Lineage Commitment}

Among the most consequential decisions a cell makes is the irreversible
commitment to a particular developmental lineage. In the blood-forming system, a
common myeloid progenitor must resolve between the erythroid fate---producing
oxygen-carrying red blood cells---and the myeloid fate, giving rise to the
granulocytes and monocytes of the innate immune system. This decision is
regulated by two cytokines: erythropoietin (EPO) promotes erythroid commitment,
and granulocyte colony-stimulating factor (GCSF) promotes myeloid commitment.
The molecular arbiters of the decision are two master transcription factors,
GATA1 and PU.1, which antagonise each other through a double-negative feedback
loop. When GATA1 dominates, erythroid genes are expressed and PU.1 activity is
suppressed; when PU.1 dominates, the myeloid programme proceeds and GATA1 is
degraded. The architecture of this toggle switch has been known since the
mid-1990s, yet a fundamental question about how it operates has resisted
resolution for more than thirty years: does the cytokine environment instruct
fate by deterministically driving the switch, or does the cell make a stochastic
choice that cytokines merely bias in probability?

\subsection{Three Decades of Unresolved Debate}

The question was first sharpened by Metcalf~\cite{metcalf1988}, who observed
from bulk colony assays that specific cytokines reliably produced specific blood
cell lineages, suggesting instruction. Cantor and Orkin~\cite{cantor2002} later
described the transcriptional architecture of erythropoiesis in mechanistic terms
that reinforced the instructive view: EPO activates GATA1 through
receptor-coupled signalling, GATA1 suppresses PU.1 through direct interaction,
and the outcome is therefore predictable from the input. Yet when single-cell and
clonal experiments became possible, the picture changed substantially. Enver
\etal~\cite{enver1998} noted that progenitors showed fate heterogeneity that
could not be explained by instructive models, asking whether stem cells play
dice. Hoppe \etal~\cite{hoppe2016} resolved individual cell fate by live imaging
of GATA1 and PU.1 reporters in clonal populations under identical cytokine
exposure and found that cell fate was determined before any EPO-dependent protein
ratio divergence at the population level---a result incompatible with pure
instruction. Olsson \etal~\cite{olsson2016} used single-cell RNA sequencing to
trace progenitor trajectories and observed that mixed-lineage states preceded
binary commitment, consistent with a system that commits stochastically from a
primed intermediate. Meanwhile, Chang \etal~\cite{chang2008} demonstrated in a
different progenitor context that transcriptome-wide noise quantitatively
predicted lineage choice, directly implicating molecular stochasticity in fate
decisions.

Both bodies of evidence are internally consistent. The apparent contradiction
between them has driven thirty years of parallel experimental traditions---bulk
assays producing instructive-looking data and single-cell assays producing
stochastic-looking data from the same biological system. A mechanistic
explanation for how both observations can simultaneously be true has not been
provided.

\subsection{Prior Computational Models and Their Limits}

The GATA1/PU.1 system has attracted substantial modelling attention. Huang
\etal~\cite{huang2007} constructed an ordinary differential equation model of the
double-negative feedback and demonstrated that the circuit produces bistability:
two stable attractors corresponding to erythroid and myeloid committed states,
separated by an unstable equilibrium. Their model captured the qualitative
phenomenology but, being deterministic, could not address fate variability at
single-cell copy numbers. Chickarmane \etal~\cite{chickarmane2009} extended this
work by introducing cooperativity in the cross-repression terms, identifying
primed-state residence time as a key predictor of commitment delay. These models
established the bistable framework within which the current debate operates, but
they address population-level attractors rather than the per-cell probability of
entering each attractor.

What has been missing is a model that explicitly treats molecular copy number as
a variable, spans the range from single-cell to population scale, and
decomposes the dynamics into mechanistic layers so that the role of noise can be
localised within the circuit. Stochastic simulation of gene regulatory networks
is well-established~\cite{gillespie1977}, but its application to the full
haematopoietic commitment circuit---including receptor kinetics, energy
metabolism, phosphorylation cascades, and the effect of nuclear pH on
transcription kinetics---has not been carried out in a unified framework.

\subsection{Approach and Contributions}

In this work we construct a stochastic model of the GATA1/PU.1 commitment
circuit using Signal Hierarchical Petri Nets~\cite{shypn2026}, a hybrid
formalism that accommodates both discrete molecular-count dynamics and continuous
biochemical kinetics within the same model structure. The modelling formalism is
described in full in the final section of this paper; here we focus on the
biological findings it enables.

The model integrates cytokine receptor binding and recycling, mRNA transcription
and nuclear export, protein translation, nuclear import, phosphorylation,
ubiquitin-mediated degradation, and the energy metabolism that powers these
transitions. Consistent with the established biology that uncommitted progenitors
carry a slight GATA1 bias before cytokine exposure~\cite{cantor2002}, we encode a
modest kinetic advantage for GATA1 transcription over PU.1 as a structural
modelling assumption. By scaling initial molecular concentrations to one hundredth
of standard values, we operate in the regime of physiological single-cell copy
numbers, where stochastic fluctuations dominate.

We performed complementary sweeps totalling more than three thousand stochastic
replicates across EPO concentration, GCSF concentration, cytokine pulse
duration, and nuclear pH. Our central findings are these. First, at single-cell
copy numbers, erythroid commitment probability is statistically flat across a
wide EPO range: stochastic noise in the decision layer, not cytokine
concentration, is the proximate determinant of individual fate. Second, the true
bifurcation is controlled by the ratio of GCSF to EPO receptor occupancy, not
by either cytokine alone, with an abrupt all-or-nothing switch between the
erythroid and myeloid attractors at a precise ratio threshold. Third,
mean-trajectory analysis reveals a three-layer cascade---reception, decision,
and execution---whose layers are EPO-dose-insensitive by distinct mechanisms.
Fourth, a cytokine-withdrawal experiment demonstrates that the erythroid state
established by EPO signalling is maintained by ongoing receptor occupancy, not
epigenetically locked; the GATA1/PU.1 crossover identified in the mean
trajectory is a cytokine-dependent priming event rather than an irreversible
molecular record. Fifth, and constituting the largest dataset reported here, a
factorial EPO--pH sweep across 2400 stochastic replicates shows that nuclear pH
shifts the EPO commitment threshold by a magnitude and in a direction that
net-flux analysis alone cannot predict, that more than half of the tested
conditions are genuinely bistable by quantitative criterion, and that a single
EPO concentration near the bifurcation can drive opposite majority fates at
acidic versus alkaline pH.

These findings reconcile the instructive and stochastic views mechanistically
and generate five experimentally testable predictions providing rational guidance
for single-cell wet-lab experiments, including a precisely bounded EPO
pulse-duration window and a direct test for the population-averaging origin of
the classical commitment threshold.
